\documentclass{article}
\usepackage[T1]{fontenc}
\usepackage{lmodern}
\usepackage{graphicx}
\usepackage{amsmath,amssymb}
\usepackage[a4paper,margin=2cm]{geometry}
\usepackage[absolute,overlay]{textpos}
\setlength{\TPHorizModule}{1mm}
\setlength{\TPVertModule}{1mm}
\usepackage[indonesian]{babel}
\usepackage{hyperref}
\setlength{\emergencystretch}{2em}
% unit: Halaman 87

\begin{document}

% === Halaman 87 (python/native) ===

87

• Contoh: 

 Plainteks: kripto (10 17  8  15  19  14) 

 n = 26, ambil m = 7   (7 relatif prima dengan 26)

 Enkripsi: C  7P + 10 (mod 26) 



p1 = 10  → c1  7  10 + 10  80  2 (mod 26) 

(huruf ‘C’)


p2 = 17  → c2  7  17 + 10  129  25 (mod 26) (huruf ‘Z’)


p3 = 8    → c3  7  8 + 10  66  14 (mod 26) 

(huruf ‘O’)


p4 = 15  → c4  7  15 + 10  115  11 (mod 26) (huruf ‘L’)


p5 = 19  → c5  7  19 + 10  143  13 (mod 26) (huruf ‘N’)


p6 = 14  → c6  7  14 + 10  108  4 (mod 26) (huruf ‘E’)

 Cipherteks: CZOLNE

\end{document}
